\subsection*{Problema 40}
\textbf{Enunciado:} Calcular la integral:
\begin{equation*}
\int x^2 (1 + x^2)^{-1/2} dx
\end{equation*}

\textbf{Solución:}
Para resolver la integral dada, utilizaremos el método de integración por partes. 
Sea la integral:
\begin{equation*}
I = \int \dfrac{x^2}{\sqrt{1 + x^2}} \, dx
\end{equation*}

Definimos las partes para aplicar la fórmula $\int u \, dv = u v - \int v \, du$:
\begin{align*}
u &= x \implies du = dx \\
dv &= \dfrac{x}{\sqrt{1 + x^2}} \, dx
\end{align*}

Para encontrar $v$, resolvemos la integral de $dv$ mediante un cambio de variable simple $w = 1 + x^2$, de donde $dw = 2x \, dx$:
\begin{align*}
v &= \int \dfrac{x}{\sqrt{1 + x^2}} \, dx \\
&= \dfrac{1}{2} \int w^{-1/2} \, dw = \sqrt{w} = \sqrt{1 + x^2}
\end{align*}

Aplicando la fórmula de integración por partes:
\begin{align*}
I &= x \sqrt{1 + x^2} \\
&\quad - \int \sqrt{1 + x^2} \, dx
\end{align*}

La integral restante se resuelve usando la sustitución trigonométrica $x = \tan\theta$, resultando en:
\begin{align*}
\int &\sqrt{1 + x^2} \, dx = \dfrac{x}{2}\sqrt{1 + x^2} \\
&+ \dfrac{1}{2}\ln\left|x + \sqrt{1 + x^2}\right|
\end{align*}

Sustituyendo este resultado de vuelta en nuestra expresión para $I$:
\begin{align*}
I &= x \sqrt{1 + x^2} - \left( \dfrac{x}{2}\sqrt{1 + x^2} \right. \\
&\quad \left. + \dfrac{1}{2}\ln\left|x + \sqrt{1 + x^2}\right| \right)
\end{align*}

Simplificando los términos semejantes, obtenemos la solución final:
\begin{align*}
I &= \dfrac{x}{2}\sqrt{1 + x^2} \\
&\quad - \dfrac{1}{2}\ln\left|x + \sqrt{1 + x^2}\right| + C
\end{align*}

\begin{figure}[H]
	\centering
	\includegraphics[width=\columnwidth]{semana_05/graficas/SEM05_P40_grafica.png}
	\caption{Representación gráfica de $f(x)$ y su antiderivada $F(x)$ evaluada para la constante de integración $C = 0$ (Problema 40).}
	\label{fig:problema40}
\end{figure}
